\documentclass{article}
\usepackage{amsmath,amssymb,mathtools}
\usepackage[margin=1in]{geometry}
\allowdisplaybreaks
\usepackage[UTF8]{ctex}
\begin{document}

\section*{选定的双层表述方式}
\begin{equation}
\begin{aligned}
\min_{(\mathbf X,\mathbf f)}\quad
\mathrm{PI}(\mathbf X,\mathbf f)
={}&\sum_{k\in K}(c_kx_k+h_ku_k)\\
&+v_1\left(\sum_{w\in W}\sum_{p\in P^w}\sum_{n\in N}t_p^wf_p^{w,n}
+\sum_{k\in K}W_k\lambda_k\right)
+v_2\sum_{a\in A}q_a^ge_a^g,\\
\text{s.t.}\quad&\mathbf X(\mathbf x,\mathbf u,\mathbf y)\in\Omega,\\
&\mathbf f\in\operatorname*{argmin}\left\{Z(\mathbf f):\mathbf f\in\Theta\right\},\\
Z(\mathbf f)={}&-\sum_{w\in W}\sum_{n\in N}q^{w,n}
\mathbb E\left[\min_{p\in P^w}\{t_p^w\}\mid\mathbf t^w(\mathbf f)\right]\\
&+\sum_{a\in A}q_at_a(q_a)-\sum_{a\in A}\int_0^{q_a}t_a(q_a)\,dq.
\end{aligned}\tag{23}
\end{equation}
这是Tran、Keyvan-Ekbatani、Ngoduy和Watling所著的运筹学教材第15页的物理PDF内容（打印版为第13页），标题为\emph{电动汽车快速充电设施部署中的随机性与环境成本纳入问题}，已接受版本（2021年）。此处显示的$Z$定义是打印版中的argmin表达式，仅为了排版而在此处分离。特别地，打印版中的被积函数$t_a(q_a)$和微分项$dq$保持不变；被积函数并未更改为$t_a(q)$。

\section*{确定可行的空间范围和随机范围}
该论文在物理页面第13页上明确定义了$\Omega$，具体内容如下：
\begin{align}
u_k&\leq mx_k &&\forall k\in K,\tag{6}\\
\mathbb G_{s,p}^{w,n}&\leq\alpha+(1-y_p^{w,n})
&&\forall w\in W,n\in N,p\in P^w,\tag{7}\\
x_k&=\{0,1\} &&\forall k\in K,\tag{8}\\
y_p^{w,n}&=\{0,1\} &&\forall w\in W,n\in N,p\in P^w,\tag{9}\\
u_k&\geq0 &&\forall k\in K.\tag{10}
\end{align}
在(7)中出现的概率，是由以下方式定义的：
\begin{equation}
\mathbb G_{s,p}^{w,n}=\mathbb P\left(D^n\leq l_{s,p}^w\right)
=G^n(l_{s,p}^w).\tag{3}
\end{equation}
这里的$\mathbb G_{s,p}^{w,n}$是一个概率值，而$G^n$则是随机驾驶范围的分布函数。在假设(I)下，对于给定的实现结果$\omega$，每个子路径的$D^n(\omega)$值都是相同的。所提供的假设表明：
\[
G^n:\Re_+^n\longrightarrow[0,1],\qquad \alpha\in[0,1].
\]
驾驶范围是指车辆在充满电的情况下能够行驶的最大距离。对于给定的充电站布局，子路径包括从起点到第一个充电站、连续充电站之间的路径，以及从最后一个充电站到目的地的路径。数值$l_{s,p}^w$表示路径$p$中，从O-D点对$w$出发的子路径$s$的长度。第12页的说明中提到，耗尽风险指的是驾驶范围小于最长子路径长度的概率。

所选的随机约束条件（7）与条件（4）重复出现。作者明确使用这一方式来替代确定性约束条件（1）和（2）。被替代的约束条件在此不再适用。阈值适用于条件（3）中的概率值，而非严格的不等式则体现在条件（7）中。

\paragraph{印刷范围与不一致性。}
方程式(8)和(9)将等式打印到$\{0,1\}$处，而不是打印出集合成员关系；这些符号被保留了下来。虽然$u_k$被称为充电器的数量，但(6)--(10)并没有为其打印出整数限制。方程式(7)中有$s$下标，但其打印出的量化因子仅涵盖$w,n,p$；这里没有缺少任何量化因子。第13页的开头段落将耗尽风险描述为“驾驶距离超过子路径长度的最小概率”，然后将其描述为小于最长子路径的概率。这些事件描述之间存在冲突。在(10)之后的解释中，成功完成的情况被描述为概率最多为$\alpha$。实际的概率仍然按照打印的定义(3)来呈现。文中将(6)称为根据预算对站点数量的限制，而其打印出的行则限制了每个节点上的充电器的数量；这里并没有添加任何预算限制。

\section*{路径-流可行空间}
物理页面第15页指出，$\Theta$是由流量守恒原理以及以下约束条件所定义的，而$M$则是一个较大的正数。
\begin{align}
f_p^{w,n}&\leq My_p^{w,n}
&&\forall w\in W,n\in N,p\in P^w,\tag{21}\\
f_p^{w,n}&\geq0
&&\forall w\in W,n\in N,p\in P^w.\tag{22}
\end{align}
流量守恒被保留为作者在$\Theta$上的设定条件。在提供的模型页面中并未显示明确的流量守恒行；其他地方的明确定义行未被验证。特别是，没有将推断的需求总和方程代入该来源参考中。

\section*{需求、路线选择以及网络映射}
在假设(II)下，车辆类别$n$在O-D对$w$之间的需求遵循独立的平稳泊松分布，其均值和方差分别为$q^{w,n}>0$。按照该假设的表示方式，
\[
Q^{w,n}\sim\operatorname{Poisson}(q^{w,n}),\qquad
\mathbb E[Q^{w,n}]=\operatorname{Var}(Q^{w,n})=q^{w,n}>0.
\]
旅行者独立地选择备选路径$p\in P^w$，每个路径的概率为$\mathbb P_p^{w,n}$，分别对应$n\in N$和$w\in W$。物理版第14页还附有相关内容。
\[
Q^w\sim\operatorname{Poisson}(q^w),\qquad
F_p^{w,n}\sim\operatorname{Poisson}(f_p^{w,n}),\qquad
\mathbf f=\mathbf f(\mathbf X).
\]
路径-流量随机变量被声明为独立的。在提供的模型描述中，并未提供$q^w$与$q^{w,n}$之间的明确关系，也没有进行任何推断。

路线选择概率在物理书的第14页上有所定义。
\begin{equation}
\mathbb P_p^{w,n}
=\mathbb P\left(t_p^w\leq t_l^w,\ \forall l\in P^w,l\ne p
\mid y_p^{w,n}=1\right).\tag{15}
\end{equation}
与其相关的平衡特性描述，应直接打印在之前的位置。所有数字都应保持不变。
\begin{equation}
f_p^{w,n}=q^{w,n}\mathbb P_p^{w,n}\left(\mathbf t^w(f)\right)
\qquad\forall w\in W,n\in N,p\in P^w.\tag{14}
\end{equation}
方程（14）描述了作者所提出的随机均衡状态，这种均衡状态可以通过方程（23）中的低层次等效最小化方法来描述。这并不是一个额外的优化问题，也不是对 argmin 施加的额外条件。所选的较低层次仍然保持为方程（23）。

链接-流映射是
\begin{align}
Q_a&=\sum_{w\in W}\sum_{p\in P^w}\sum_{n\in N}
\delta_{a,p}^wF_p^{w,n} &&\forall a\in A,\tag{16}\\
q_a&=\sum_{w\in W}\sum_{p\in P^w}\sum_{n\in N}
\delta_{a,p}^wf_p^{w,n} &&\forall a\in A.\tag{17}
\end{align}
链接流随机变量$Q_a$被描述为具有边际泊松分布，其均值为$q_a$。

\section*{行程时间函数}
在物理页面14上，未编号的多项式关系以及选定的二次阶数，应保持不变。所有数字都应原封不动地保留下来。
\[
t_a(q_a)=\sum_{j=0}^{e}b_{ja}q_a^j,\qquad e=2.
\]
作者采用了 BPR 函数，并指出使用泰勒级数近似法来计算$b_{ja}$的值。他们展示的旅行时间族为：
\begin{align}
t_a&=t_a^0\left[1+0.15\left(\frac{q_a}{C_a}\right)^4\right]
&&\forall a\in A,\tag{18}\\
\hat t_a&=t_a+b_{2a}q_a &&\forall a\in A,\tag{19}\\
t_p^w&=\sum_{a\in A}\hat t_a\delta_{a,p}^w
&&\forall p\in P^w.\tag{20}
\end{align}
据称，所有车辆类别的感知路径行驶时间都是相同的，其平均值由公式(20)给出。表2中列出了相关数据。
\[
\mathbf t^w=(t_p^w)^T.
\]
系数$b_{ja}$仍然保持抽象状态：所提供的页面并未指定泰勒展开点或具体的系数公式。多项式近似和BPR表达式则按照作者所描述的角色进行保留，而不是将(18)中的四次幂替换为二次幂。

\paragraph{按原文格式标注。}
表2中将$t_p^w$定义为随机路径行驶时间$T_p^w$的平均值。而在(15)和(23)中的期望值计算中，概率表达式内使用了相同的小写字母$t_p^w$。这些打印出的符号被保留了下来。关于感知误差的分布情况，实际上在附录A的最后一段中有详细说明，具体位置在物理PDF页面的24页（打印页数为22页）。该内容被包含在下面；这并不是一个无法获取的外部定义。在表2中，容量为$c_a$，而在(18)中则为$C_a$。显示的大写符号被保留下来，同时没有在这两个符号之间添加等号。

\paragraph{附录A中列出的感知错误分布。}
链接感知错误在各个链接之间是独立的，并且对所有车辆类别来说，其分布规律都是相同的。对于其他未命名的链接错误，我们使用辅助符号$\varepsilon_a$来表示。打印出的正常法则如下：
\[
\varepsilon_a\sim\mathcal N\!\left(0,\bigl(\phi t_a(0)\bigr)^2\right),
\qquad a\in A,
\qquad \{\varepsilon_a:a\in A\}\text{ independent}.
\]
这里，$t_a(0)$表示零流量链接行驶时间，而$\phi$表示感知误差尺度。该论文在数值测试中使用了$\phi=0.3$的值；该值并未被应用于一般公式。其中“所有车辆类别均相同”描述了共同的链接误差规律；并未断言车辆类别之间或重叠路径之间存在额外的独立性。这一规律，连同经过泊松修正的链接行驶时间和路径发生率，构成了在(14)--(15)中命名的probit感知模型。在(15)和(23)中，使用小写路径时间符号的印刷方式并未通过添加不同的公式来隐式修正。

假设（III），物理页面10（打印页面8）指出，旅行者在选择路线时，会考虑路径的可行性和连接点的旅行时间。因为假设充电站的等待时间对选择路线来说并不重要。因此，充电站的等待时间仍然包含在高层成本中，但不会被计入公式（20）中计算出的路径时间。

\section*{拥堵收费}
假设和上层定义中未编号的功能依赖关系是
\[
\lambda=\lambda(\mathbf f),\qquad
W_k=W_k(u_k,\lambda_k,\mu_k).
\]
据称，充电站的到达率是根据相应的 SUE 流量来推断出来的。每个充电器的充电时间都是独立且相同分布的指数型时间，在假设中表示为 $\mu$。对于固定分布的情况，到达率被假设为严格小于服务率。这一严格假设被保留下来，没有被削弱为非严格的比较，也没有被替换为虚构的站点索引不等式。

在物理页面14上进行的递归期望时间计算是：
\begin{align}
a_1&=1;\qquad a_i=1+\frac{\mu}{\lambda}(i-1)a_{i-1},\tag{11}\\
W(u,\lambda,\mu)&=
\frac{\lambda}{(u\mu-\lambda)^2
\left[a_u+\dfrac{\lambda}{u\mu-\lambda}\right]}+\frac{1}{\mu}.\tag{12}
\end{align}
这里的$u$、$\lambda$和$\mu$是等待时间函数的站点级参数，而$a_i$则是递归量，与链接索引$a$不同。在(11)旁边没有打印出$i$的范围。物理页面13中的定义将$\mu_k$称为站点服务率，而表2则将其称为收费时间；(11)--(12)及其服务率的使用方式保持不变，没有引入任何新的递归定义。表2中使用的$w_k$是小写字母形式，表示预期等待时间；而目标和功能定义中使用的$W_k$则是大写字母形式。

\paragraph{从附录B中得到的到达率计算公式。}
附录B中，物理PDF页面24-25（印刷版为22-23页）详细说明了EV路径如何决定车辆到达各个站点的情况。其中，$l$被定义为$\mathbb P(D^n\leq l)$小于$\alpha$的最大长度。对于具有安装充电站的路径，需要确定从哪些之前的充电点可以到达每个后续的充电点或目的地。旅行者会在这些可行的充电点之间平均分配。沿着路径反向进行同样的推理，将每个站点需要重新充电的所有情况加总起来。总的站点到达率就是各个EV路径到达率的总和。仅仅通过某个站点并不意味着需要在该站点进行充电，因此，总的站点到达率并不是所有经过该站点的路径流量的总和。

为了精确起见，以下未编号的记账条目重新阐述了附录中的逆向等比例构建方式；其中的辅助符号属于编辑性标注，并非额外的规划决策或新分配的等式编号。固定一个EV路径及其类别，该路径的流量为$f_p$。按路径顺序列出起点、已安装的充电站以及终点。对于每一个非起点的位置$j$，让$B(j)$成为其可行的前驱：即更早的充电位置，或者满电状态下的起点，从起点到$j$的沿路径距离最多为附录中的范围$l$。将终点所需的到达质量设定为$r_d=f_p$，并将所有更早的质量初始化为零。按逆向路径顺序访问终点以及各个充电位置。在每个位置$j$处，分配其累积的质量。
\[
 r_i\ \leftarrow\ r_i+\frac{r_j}{|B(j)|}
 \qquad \text{for each }i\in B(j).
\]
起点作为起点，并不以充电到达点身份出现。路径对充电事件的贡献$j$即是其累积值$r_j$。若一个节点存在多个充电事件，则各事件的贡献值相加。将这一构造应用于每个具有自身范围和流量的电动汽车类别，然后对路径、O-D对及电动汽车类别进行求和，最终得到$\lambda_k$。这一最终求和即为附录中所述的聚合结果，并非假设每个访问节点$k$的电动汽车都会在那里充电。该构造需要可行的充电链；论文并未在所需的前置集为空时定义任意到达率。

打印示例中的节点序列为$1,2,3,4,5,6$，充电站为$2,3,4$，连续链接的长度分别为$25,30,10,10,25$公里，路径流量为$f_p=100$辆/小时和$l=50$公里。目的地$6$可以在$3$或$4$充电后到达，两种充电方式的比例相同。充电站$4$则可以在$2$或$3$充电后到达，同样两种充电方式的比例相同；而充电站$3$则需要在$2$充电后才能到达。
\[
\lambda_{4,p}=\tfrac12(100)=50,\qquad
\lambda_{3,p}=\tfrac12(100)+\tfrac12(50)=75,\qquad
\lambda_{2,p}=75+\tfrac12(50)=100.
\]
这些是图 B.11 中打印出的费率。这个例子说明了一般规律；其数值流程、链接长度以及范围并不是每个实例的固定输入。

\paragraph{保留源边界问题。}
附录B中提到了“小于”的$\alpha$，而操作约束条件（7）则使用了$\leq\alpha$。此外，附录B还将$l$称为最大值，但没有对未达到严格阈值的端点进行任何规定。两种来源的表述都被记录下来；没有发明任何逆CDF端点规则来调和它们。方程（11）到（12）中并没有给出$\lambda=0$或$u=0$的值；对于$u_k$也没有提供任何限制值或正下限的规定。

\section*{汽油车排放}
排放函数及其固定系数如下：
\begin{equation}
e_a^g=\rho_1\hat t_a\exp\left(\rho_2\frac{l_a}{\hat t_a}\right).\tag{13}
\end{equation}
\[
\rho_1=11.063927,\qquad \rho_2=0.008493.
\]
量$e_a^g$是汽油车辆在链接$a$上平均排放的CO量。目标函数使用了它们的平均链接流量$q_a^g$和CO单位成本$v_2$。在这项道路计算中，电动车被视为零排放车辆。表2将$q_a^n$定义为车辆类别$n$的平均链接流量；在提供的模型页面中并未显示针对$q_a^g$的单独类别特定路径到链接的总和。这里没有用推断的公式或假设的$N$枚举来替代它。

\section*{集合、固定的输入以及可变的角色}
以下是表2（物理页数9-10）中的操作含义以及相关的模型文本。
\begin{description}
\item[网络集合] $K$是节点集合，其中$k\in K$为节点；$A$是链接集合，其中$a\in A$为链接；$W$是O-D对集合，其中$w\in W$为O-D对；$N$是车辆类别集合，其中$n\in N$为车辆类别；$P^w$是O-D对$w$之间的路径集合，其中$p\in P^w$为路径。这些符号表示集合，而不是它们的基数。没有规定具体的车辆类别数量或枚举方式。
\item[网络输入] $t_a^0$表示自由流动链接的传输时间；$c_a$表示表2中的链接容量；$l_a$表示链接长度；$l_p^w$表示路径长度；$l_{s,p}^w$表示子路径长度。这些数值是根据特定条件定义的。
\[
\delta_{k,p}^w=
\begin{cases}1,&\text{node }k\text{ is on path }p\text{ between pair }w,\\0,&\text{otherwise},\end{cases}
\qquad
\delta_{a,p}^w=
\begin{cases}1,&\text{link }a\text{ is on path }p\text{ between pair }w,\\0,&\text{otherwise}.\end{cases}
\]
\item[规划与成本输入] $m$表示每个站点可容纳的最大充电设备数量；$c_k$表示在节点$k$开设一个站点的固定成本；$h_k$表示在该地点安装一个充电设备的单位成本；$v_1$表示旅行时间的货币价值；$v_2$表示每吨二氧化碳的单位成本。这些都是固定数据，并非需要优化的变量。需求、行驶里程分布、风险阈值以及充电参数都受到上述限制；没有进一步的参数限制信息提供。
\item[规划决策]。$x_k$表示某个站点是否位于$k$位置；$u_k$是放置在该站点的充电器数量；$y_p^{w,n}$表示路径$p$在$w$对对于$n$类别是否可行。它们的向量分别为$\mathbf x,\mathbf u,\mathbf y$和$\mathbf X$，后者负责收集这三个向量。表2将路径可行性符号标记为$y_{p,n}^w$，而操作方程则使用$y_p^{w,n}$。
\item[平衡量.] $F_p^{w,n}$是随机的类别特定路径流量，$f_p^{w,n}$是其平均值，而$\mathbf f$是所有平均路径流量的向量。$Q_a$是随机的聚合链接流量，$q_a$是其平均值。$T_a$是随机的链接旅行时间，而$t_a$在表2中是其平均值；$T_p^w$是随机的路径旅行时间，$t_p^w$是其平均值。选定的公式通过较低层次的平衡和显示的网络函数来确定这些量，而不是将所有这些量视为独立的规划决策。
\item[充电量] $\lambda_k$是车站$k$处的平均EV到达率，该数值取决于均衡流量。$W_k$则是目标函数中预期使用的時間，该数值是通过车站函数以及(11)--(12)公式来估算得出的。
\end{description}

\paragraph{验证边界。}
所选公式及其相关定义已与原始PDF文件进行核对，包括第13-15页的物理公式、第9-11页的符号和假设，以及第23-25页的附录A/B内容。上述两个附录定义均来自同一份PDF文件，而非独立的未验证补充材料。源文件中的编号数学表达式被保留下来；新增的正常律符号、反向到达计算以及算术示例被明确标示为对应定义的具体体现。第17页的数值网络假设，包括Gamma驱动范围以及风险阈值$0.05$，属于案例研究中的选择，并非对(23)的强制性要求。同样，附录A中的实验值$\phi=0.3$被明确标注，但并未固定通用参数。CEM和MSA求解算法并未被替代为所选的较低级优化方法。关于等效优化、行程时间近似和排队公式的外部参考并未单独提供，且超出此PDF文件中打印的表达式和定义部分仍未被验证。保留并注释源文件中的不一致性，体现了对所选公式的忠实再现，但这并不意味着这些源文件中的不一致性在数学上得到了解决。


\paragraph{离散充电领域的额外原始证据。}
不等式 (6) 本身并未明确说明积分性。在附录 A 中，(A.1) 通过节点上的计数来索引概率矩阵，计数范围从零到 \(m\)。在物理 PDF 的第 24 页上，作者对所有 \(\mathbf X\in\Omega\) 在 (A.2) 中进行最小化，并将某个节点的干预行为指定为 \(X_i=j\)，其概率为 \(\gamma_{i,j}\)，对应的数量为 \(j=1,\ldots,m\)。一个封闭的节点没有充电点。这些基准条件要求 EV 需求为正，从而迫使每个候选节点通过旅行概率约束 (3),(7) 保持开放状态。队列递归运算应用于这些整数计数。每个基准条件都拥有自己的泰勒中心；唯一的路径及其严格划分的可行与不可行范围比较，使得可以在有限的实际子路径距离上进行直接的概率检验。评估这些有限路径时，不需要对开放均匀 CDF 子层次集合进行连续最大值分析，也不假设在分数充电点存在队列值。
\end{document}